%! AP Statistics — Unit 2, Lesson 2.7 — Slides
\documentclass[aspectratio=169, 11pt]{beamer}
\usepackage{apstats-beamer}

\begin{document}

% ── TITLE SLIDE ───────────────────────────────────────────────────────────────
{
\setbeamercolor{background canvas}{bg=navy}
\begin{frame}[plain]
  \vspace{2.8cm}\hspace{0.55cm}%
  \begin{minipage}{0.9\textwidth}
    {\color{goldacc}\bfseries\small LESSON 2.7}\\[0.5cm]
    {\color{white}\bfseries\LARGE Residuals}\\[1.2cm]
    {\color{skymid}\small AP Statistics~$\cdot$~Unit 2: Exploring Two-Variable Data}
  \end{minipage}
\end{frame}
}

% ── LEARNING TARGETS ─────────────────────────────────────────────────────────
\begin{frame}[t, plain]
  \navyheader{Learning Targets}
  \begin{enumerate}
    \item I can calculate a residual and interpret it in context as a measure of prediction error.
          \hfill\textcolor{slate}{\small(DAT-1.E)}
    \bigskip
    \item I can construct a residual plot by graphing residuals vs.\ the explanatory variable.
          \hfill\textcolor{slate}{\small(DAT-1.E)}
    \bigskip
    \item I can describe the pattern in a residual plot and assess whether a linear model is
          appropriate for the data.
          \hfill\textcolor{slate}{\small(DAT-1.F)}
  \end{enumerate}
\end{frame}

% ── THE RESIDUAL ─────────────────────────────────────────────────────────────
\begin{frame}[t, plain]
  \navyheader{What Is a Residual?}
  \sectionlabel[goldacc]{DAT-1.E}

  \begin{columns}[T]
    \column{0.50\textwidth}
      \textbf{The formula:}
      \[
        e = y - \hat{y} \quad \text{(actual $-$ predicted)}
      \]
      \bigskip
      \begin{tabular}{ll}
        $e > 0$: & point is \textbf{above} the line \\[4pt]
        $e < 0$: & point is \textbf{below} the line \\[4pt]
        $\sum e_i = 0$ & always true for the LSRL
      \end{tabular}

    \column{0.46\textwidth}
      \textbf{Example}\\[4pt]
      $\widehat{\text{temp}} = 36 + 3(\text{sunlight})$\\[6pt]
      City with 6 hr sunlight, actual temp = 52°F:
      \[
        \hat{y} = 36 + 3(6) = 54^\circ\text{F}
      \]
      \[
        e = 52 - 54 = -2^\circ\text{F}
      \]
      {\small The actual temp is 2°F \textbf{below} the prediction.}
  \end{columns}
\end{frame}

% ── RESIDUAL TABLE ───────────────────────────────────────────────────────────
\begin{frame}[t, plain]
  \navyheader{Computing All Residuals}
  \sectionlabel[goldacc]{DAT-1.E}

  \[
    \widehat{\text{temp}} = 36 + 3(\text{sunlight})
  \]

  \vspace{0.1in}
  \renewcommand{\arraystretch}{1.6}
  \begin{tabular}{cccr}
    \rowcolor{sky}
    \textbf{Sunlight ($x$)} & \textbf{Temp ($y$)} & \textbf{Predicted ($\hat{y}$)} &
    \textbf{Residual ($e$)} \\
    6  & 52 & 54 & $-2$ \\
    8  & 62 & 60 & $+2$ \\
    10 & 66 & 66 & $0$ \\
    12 & 74 & 72 & $+2$ \\
    14 & 76 & 78 & $-2$ \\
    \hline
    \multicolumn{3}{r}{\textbf{Sum}} & $\mathbf{0}$ \\
  \end{tabular}

  \vspace{0.15in}
  {\small Check: residuals always sum to \textbf{zero} for the least-squares line.}
\end{frame}

% ── RESIDUAL PLOT ────────────────────────────────────────────────────────────
\begin{frame}[t, plain]
  \navyheader{Constructing a Residual Plot}
  \sectionlabel[goldacc]{DAT-1.E}

  \begin{columns}[T]
    \column{0.48\textwidth}
      \textbf{How to build one:}
      \begin{enumerate}
        \item New axes: $x$ = explanatory variable, $y$ = residual
        \item Plot each $(x_i,\; e_i)$
        \item Draw a dashed line at $e = 0$
      \end{enumerate}
      \bigskip
      \textbf{What to look for:}
      \begin{itemize}
        \item \textcolor{greenacc}{\textbf{Random scatter}} $\to$ linear OK
        \item \textcolor{redacc}{\textbf{Curved pattern}} $\to$ linear NOT appropriate
        \item \textcolor{redacc}{\textbf{Fanning}} $\to$ investigate further
      \end{itemize}

    \column{0.48\textwidth}
      \begin{tikzpicture}[scale=0.9]
        \draw[->] (0,0) -- (7,0) node[right]{\small sunlight (hr)};
        \draw[->] (0,-1.4) -- (0,1.4) node[above]{\small residual (°F)};
        \draw[dashed,gray] (0,0) -- (7,0);
        \foreach \x/\lbl in {1.2/6, 2.4/8, 3.6/10, 4.8/12, 6.0/14}
          \draw (\x,0.07)--(\x,-0.07) node[below]{\tiny\lbl};
        \foreach \y/\lbl in {-1.2/{$-2$}, 1.2/{$+2$}}
          \draw (0.07,\y)--(-0.07,\y) node[left]{\tiny\lbl};
        \filldraw[navy] (1.2,-1.2) circle (3pt);
        \filldraw[navy] (2.4, 1.2) circle (3pt);
        \filldraw[navy] (3.6,   0) circle (3pt);
        \filldraw[navy] (4.8, 1.2) circle (3pt);
        \filldraw[navy] (6.0,-1.2) circle (3pt);
      \end{tikzpicture}
      \vspace{4pt}
      {\small Random scatter $\to$ linear model appears appropriate.}
  \end{columns}
\end{frame}

% ── AP TEMPLATE ──────────────────────────────────────────────────────────────
\begin{frame}[t, plain]
  \navyheader{AP Response Template}
  \sectionlabel[goldacc]{DAT-1.F}

  \textbf{Three parts every time:}
  \begin{enumerate}
    \item \textbf{Pattern}: random scatter / curved / fanning
    \item \textbf{Conclusion}: is or is not an appropriate linear model
    \item \textbf{Context}: name the variables
  \end{enumerate}

  \bigskip

  \begin{block}{Curved pattern example}
    ``The residual plot shows a \textbf{curved pattern}. This suggests that a linear model
    \textbf{may not be appropriate} for predicting \textit{[response]} from
    \textit{[explanatory]}. A nonlinear model may provide a better fit.''
  \end{block}

  \medskip

  \begin{block}{Random scatter example}
    ``The residual plot shows \textbf{no apparent pattern} (random scatter around zero).
    This suggests that a linear model \textbf{is appropriate} for predicting
    \textit{[response]} from \textit{[explanatory]}.''
  \end{block}
\end{frame}

% ── PREVIEW ──────────────────────────────────────────────────────────────────
\begin{frame}[t, plain]
  \navyheader{Coming Up --- Lesson 2.8}

  \begin{itemize}
    \item The LSRL minimizes $\displaystyle\sum e_i^2$ (sum of \textbf{squared} residuals).
    \bigskip
    \item Why squared? Because $\sum e_i = 0$ always, so squaring is needed to penalize
          large errors.
    \bigskip
    \item Formulas: $b = r \cdot \dfrac{s_y}{s_x}$ \quad and \quad $a = \bar{y} - b\bar{x}$
    \bigskip
    \item We'll see how $r$, $s_x$, $s_y$, $\bar{x}$, and $\bar{y}$ work together to define
          the best-fit line.
  \end{itemize}
\end{frame}

\end{document}
